\chapter{Computational Hardness and One-way Functions}

Shannon secrecy is too hard to achieve in practice (a key longer than the message is not desirable). Now we relax those limits by restricting adversary ``feasible'', or ``efficient'' computation.

\section{Computational Hardness}

Here we introduce some conventions and paradigms:

\begin{definition}[Feasible / Efficient Computation]
	A computation is said to be \textbf{feasible} or \textbf{efficient} if it can be done in polynomial time, i.e. $O(n^c)$ for a constant $c$.
	$n$ is called the \textbf{security parameter}.
\end{definition}

\begin{definition}[Tiny / Negligible Function]
	A function $v(n)$ is said to be \textbf{tiny} or \textbf{negligible}, written as $v(n) = \negl(n)$, if $v(n) = O(1/n^c)$ for every constant $c$.
\end{definition}

For example, $1/2^n$ and $1/n^{\log n}$ are negligible, but $1/n^5$ is not.

\section{One-way Functions}

Intuitively, a one-way function is a function that is easy to compute but hard to invert. To formalize this, we'll do a one-wayness experiment:

\procedureblock{One-wayness ``game''}{
	\textbf{challenger} \< \< \textbf{adversary}\\
	x \gets \{0,1\}^n \< \< \\
	y = F(x) = f(x) \< \< \\
	\< \sendmessageright*{y} \< \\
	\< \sendmessageleft*{x^*\in \bit^n} \< \\
	\text{Check if }f(x^*) = y \< \<
}

\begin{remark}
	We don't simply check $x^* = x$ because there're collisions. Consider $f(x) = 0, \forall x\in\bit^n$. The probability of $x=x^*$ will have probability at most $1/2^n$. It's negligible, but such a function can't be considered secure.
\end{remark}

\begin{definition}[One-way Function (OWF)]
	A function $f: \{0,1\}^* \to \{0,1\}^*$ is called a \textbf{one-way function} if:
	\begin{itemize}
		\item Easy to evaluate: there exists a polynomial-time algorithm $F$ that evaluates $f$: $F(x) = f(x) \forall x\in \{0,1\}^*$.
		\item For any probabilistic polynomial-time (PPT) adversary $\mathcal I$,
		      \[
			      \adv_f^{\text{ow}}(\mathcal I) := \Pr_{x\gets\bit^n}[\mathcal I(1^n, y = f(x))\in f^{-1}(y)] = \negl(n)
		      \]

	\end{itemize}
\end{definition}
\begin{remark}
	$1^n$ here gives a guarantee on the input size. Only by $1^n$ can we be sure that the adversary is poly-time in $n$ instead of something else. It's just a convention and we'll omit such notations in the future.
\end{remark}

Nowadays, scientists are working on PRGs and PRFs. They comply secure encryption with $|K| \ll |M|$, many-message security, and so on.

\section{Does OWF Exist?}

The bad news is that the existence of OWFs can indicate $\P \neq \NP$. As this question remains open by the time we take the course, we don't know about the existence of OWFs.

But the good news is they appear to be abundant in practice. Here we give some examples:

\begin{eg}[Subset Sum]
	$f_{ss}: (\mathbb Z_n)^n\times \{0,1\}^n \to (\mathbb Z_n)^n\times \ZZ_n$ with $N = 2^n$.

	\[ f_{ss}(\vector{a}, \vector{b}) = (\vector{a}, s=\sum_{i=1}^n a_i b_i \bmod N) = (\vector{a}, \langle \vector{a}, \vector{b} \rangle \bmod N) \]

	The inversion game is as follows. We choose $\vector{a}\in \ZZ_N^n, \vector{b}\gets \bit^n$, and give $y = (\vector{a}, s)$ to $\mathcal I$.
	$\mathcal I$ outputs $(\vector{a}^*, \vector{b}^*)$ and wins if $f_{ss}(\vector{a}^*,\vector{b}^*) = y \iff \vector{a}^* = \vector{a}$ and $\langle \vector{a}^*, \vector{b}^* \rangle = s\bmod N$.
\end{eg}
